\clearpage
\phantomsection
\addcontentsline{toc}{chapter}{Danh mục ký hiệu và chữ viết tắt}
\chapter*{Danh mục ký hiệu và chữ viết tắt}
{\renewcommand{\arraystretch}{1.4}
{\fontsize{12}{13}\selectfont
\begin{longtable}{|C{0.8cm}|>{\raggedright}p{3.6cm}|p{9.8cm}|}
\hline
\multicolumn{3}{|l|}{\textbf{Danh mục ký hiệu}}\\
\hline\hline
\textbf{STT} & \textbf{Ký hiệu} & ~\hfill\textbf{Giải thích}\hfill~\\
\hline
1&A&Người nói của clip III\\
\hline
2&B&Người phản hồi, tức người nói của clip IV (cần dự đoán cảm xúc)\\
\hline
3&L&Người nghe được chọn làm đại diện cho B (listener proxy)\\
\hline
4&O&Những người còn lại trong cảnh\\
\hline
5&$k$&Chỉ số clip, $k \in \{1,2,3\}$ ứng với clip I, II, III\\
\hline
6&$x_{r,k}$&Token khuôn mặt của vai trò $r$ trong clip $k$\\
\hline
7&$s_k$, $g_k$&Token lời nói và token bối cảnh của clip $k$\\
\hline
8&$q$&Token truy vấn dùng để đọc kết quả dự đoán\\
\hline
9&$\phi(f)$&Vector đặc trưng 146 chiều của một khuôn mặt $f$\\
\hline
10&$p(y_B \mid X)$&Phân phối dự đoán cảm xúc của B khi biết ba clip đầu vào\\
\hline
11&$d$&Số chiều ẩn của mô hình ($d = 128$)\\
\hline
\end{longtable}
}}
{\renewcommand{\arraystretch}{1.2}
{\fontsize{12}{13}\selectfont
\begin{longtable}{|C{0.8cm}|p{1.7cm}|>{\raggedright}p{5.5cm}|p{5.75cm}|}
\hline
\multicolumn{4}{|l|}{\textbf{Danh mục chữ viết tắt}}\\
\hline\hline
\textbf{STT} & \textbf{Viết tắt} & \textbf{Giải thích tiếng Anh} & \textbf{Giải thích tiếng Việt}\\
\hline
1&CI&Confidence Interval&Khoảng tin cậy\\
\hline
2&CLIP&Contrastive Language--Image Pre-training&Mô hình học tương phản ảnh -- văn bản\\
\hline
3&CV&Cross-Validation&Kiểm định chéo\\
\hline
4&ECAPA&Emphasized Channel Attention, Propagation and Aggregation&Mạng nhúng giọng nói ECAPA-TDNN\\
\hline
5&EMAP&Empirical Multimodally-Additive Projection&Phép chiếu cộng tính thực nghiệm đa phương thức\\
\hline
6&ERC&Emotion Recognition in Conversation&Nhận diện cảm xúc trong hội thoại\\
\hline
7&FFN&Feed-Forward Network&Mạng truyền thẳng\\
\hline
8&Hi-EF&Hi-EF benchmark&Bộ dữ liệu dự đoán cảm xúc trong tương tác đa lớp\\
\hline
9&LA&Logit Adjustment&Hiệu chỉnh logit theo tần suất lớp\\
\hline
10&LN&Layer Normalization&Chuẩn hóa theo lớp\\
\hline
11&MCIS&Multilayered-Contextual Interaction Sample&Mẫu tương tác ngữ cảnh đa lớp\\
\hline
12&MELD&Multimodal EmotionLines Dataset&Bộ dữ liệu cảm xúc đa phương thức MELD\\
\hline
13&MHA&Multi-Head Attention&Chú ý đa đầu\\
\hline
14&MLP&Multi-Layer Perceptron&Mạng nơ-ron nhiều lớp\\
\hline
15&NLL&Negative Log-Likelihood&Hàm hợp lý âm\\
\hline
16&OOF&Out-of-Fold&Dự đoán ngoài fold\\
\hline
17&PCA&Principal Component Analysis&Phân tích thành phần chính\\
\hline
18&PMA&Pooling by Multihead Attention&Gộp bằng chú ý đa đầu\\
\hline
19&SAB&Set Attention Block&Khối chú ý trên tập hợp\\
\hline
20&UAR&Unweighted Average Recall&Độ nhạy trung bình không trọng số\\
\hline
21&WAR&Weighted Average Recall&Độ chính xác (độ nhạy trung bình có trọng số)\\
\hline
\end{longtable}
}}
